% Standalone proof, also input by correlation-gap-n4-full.tex.
% Candidate resolution pending expert review; frozen campaign date 2026-07-21.
\ifdefined\CGFullPaper
\let\ProofOnlyEndDocument\relax
\else
\documentclass[11pt]{article}
\usepackage[T1]{fontenc}
\usepackage{amsmath,amssymb,amsthm,mathtools}
\usepackage[margin=1in]{geometry}
\usepackage{xcolor}
\usepackage{booktabs,array}
\usepackage[hidelinks]{hyperref}
\usepackage{microtype}
\allowdisplaybreaks

\newtheorem{theorem}{Theorem}[section]
\newtheorem{lemma}[theorem]{Lemma}
\newtheorem{proposition}[theorem]{Proposition}
\theoremstyle{remark}
\newtheorem{remark}[theorem]{Remark}

\newcommand{\E}{\mathbb E}
\newcommand{\one}{\mathbf 1}
\newcommand{\D}{\mathcal D}
\newcommand{\DPI}{\mathcal D_{\mathrm{PI}}}
\newcommand{\Cfour}{C_4}
\newcommand{\Fslice}{\mathcal F}
\newcommand{\CGfour}{\mathrm{CG4}}
\newcommand{\Rge}{\mathbb R_{\ge 0}}

\title{Proof of the Four-Element Pairwise-Independent Correlation Gap}
\author{Max Laval-Mullen\\\small Independent researcher}
\date{July 22, 2026}
\begin{document}
\def\ProofOnlyEndDocument{\end{document}}
\maketitle

\begin{abstract}
We prove, subject to an exact finite certificate described below, that for
every nonnegative monotone submodular function on a four-element ground set
and every real vector of singleton marginals, the unrestricted concave
closure is at most \(4/3\) times its pairwise-independent analogue.  The
constant is sharp.  The finite premise is checked with exact rational
arithmetic and replayed by a separate Node.js/\texttt{BigInt}
implementation.
\end{abstract}

\section{Proof of the four-element theorem}\label{sec:proof}

\subsection{Statement and exact negation}\label{subsec:statement}

Let \(E=[4]\).  For a probability law \(p\) on \(2^E\), write
\(p_S=\Pr_p(S)\).  Given \(x\in[0,1]^4\), define
\[
\D(x)=\left\{p\in\Rge^{\,2^E}:\sum_{S\subseteq E}p_S=1,
\quad \sum_{S\ni i}p_S=x_i\quad(i\in E)\right\}.
\]
Let \(\DPI(x)\) be the subset additionally satisfying
\[
\sum_{S\supseteq\{i,j\}}p_S=x_i x_j\qquad(i<j).
\]
For a nonnegative monotone submodular set function
\(f:2^E\to\Rge\), put
\[
f^+(x)=\max_{p\in\D(x)}\E_p f,
\qquad
f^{\mathrm{PI}}(x)=\max_{q\in\DPI(x)}\E_q f.
\]

\begin{theorem}[Four-element pairwise-independent correlation gap]
\label{thm:main}
For every nonnegative monotone submodular
\(f:2^{[4]}\to\Rge\) and every real \(x\in[0,1]^4\),
\[
3f^+(x)\le 4f^{\mathrm{PI}}(x).
\tag{CG4}\label{eq:CG4}
\]
The factor \(4/3\) is best possible.
\end{theorem}

The exact negation of \eqref{eq:CG4} is the existence of admissible
\((f,x)\) and \(p\in\D(x)\) such that
\[
3\E_p f>4\E_q f
\qquad\text{for every }q\in\DPI(x).
\tag{NEG}\label{eq:neg}
\]
The universal quantifier over \(q\) is essential.

\subsection{Normalization}\label{subsec:normalization}

\begin{lemma}[Normalization]\label{lem:normalization}
It is enough, in both the proof and disproof directions, to consider
functions satisfying \(f(\varnothing)=0\).
\end{lemma}

\begin{proof}
Put \(c=f(\varnothing)\ge0\) and \(g=f-c\).  Monotonicity gives \(g\ge0\),
and subtracting a constant preserves monotonicity and submodularity.  Every
expectation shifts by \(c\), so
\[
f^+=c+g^+,
\qquad
f^{\mathrm{PI}}=c+g^{\mathrm{PI}}.
\tag{2.1}\label{eq:shift}
\]
If \eqref{eq:CG4} holds for \(g\), then
\[
c+g^+\le c+\frac43g^{\mathrm{PI}}
\le\frac43(c+g^{\mathrm{PI}}).
\]
Conversely, a strict violation by \(f\) gives
\[
g^+>\frac43g^{\mathrm{PI}}+\frac c3
\ge\frac43g^{\mathrm{PI}}.
\]
Thus normalization loses neither direction.
\end{proof}

Let \(\Cfour\) be the cone of normalized monotone submodular functions on
\(E\).  If \(g\in\Cfour\) and \(g(E)=0\), monotonicity forces \(g=0\).
Every nonzero member can therefore be scaled uniquely into
\[
\Fslice=\{g\in\Cfour:g(E)=1\}.
\]
This is a nonempty compact convex polytope, since
\(0\le g(S)\le1\) for every \(S\subseteq E\).

\subsection{The simultaneous-dominance reduction}
\label{subsec:minimax}

\begin{lemma}[Simultaneous dominance]\label{lem:minimax}
For fixed \(x\), \eqref{eq:CG4} for all normalized functions is equivalent
to the following statement: for every \(p\in\D(x)\), there is one
\(q\in\DPI(x)\) such that
\[
4\E_q g\ge3\E_p g\qquad(g\in\Cfour).
\tag{3.1}\label{eq:simultaneous}
\]
\end{lemma}

\begin{proof}
Fix \(p\), let \(Q=\DPI(x)\), and define
\[
L(g,q)=3\E_p g-4\E_q g
\qquad(g\in\Fslice,\ q\in Q).
\]
The product law lies in \(Q\), so \(Q\) is nonempty; it is also compact and
convex.  Finite-dimensional bilinear minimax gives
\[
\max_{g\in\Fslice}\min_{q\in Q}L(g,q)
=
\min_{q\in Q}\max_{g\in\Fslice}L(g,q).
\tag{3.2}\label{eq:minimax}
\]
For fixed \(g\), the inner minimum on the left is
\[
3\E_p g-4g^{\mathrm{PI}}(x).
\]
If \eqref{eq:CG4} holds, this is nonpositive for every \(p\) and \(g\),
because \(\E_p g\le g^+(x)\).  Equation \eqref{eq:minimax} then supplies a
single \(q\) for which \(L(g,q)\le0\) for every \(g\in\Fslice\).
Homogeneity extends the result from \(\Fslice\) to \(\Cfour\), including the
zero function.

Conversely, apply \eqref{eq:simultaneous} to a maximizer \(p\) of
\(g^+(x)\).  Then
\[
3g^+(x)\le4\E_q g\le4g^{\mathrm{PI}}(x).
\]
\end{proof}

This is the reason that one may use a common finite ray test.  Choosing a
different \(q\) for each ray would not prove
\eqref{eq:simultaneous}.

\subsection{Extreme unrestricted laws and sorting}
\label{subsec:extreme-laws}

Let \(A\) be the \(5\times16\) first-moment matrix with columns
\[
a(S)=(1,\one_{1\in S},\ldots,\one_{4\in S})^T.
\]
It has rank five.
Indeed, the columns indexed by \(\varnothing,\{1\},\ldots,\{4\}\) are
linearly independent.

\begin{lemma}[Extreme-law reduction]\label{lem:extreme-laws}
It is enough to prove \eqref{eq:simultaneous} for extreme points of
\(\D(x)\).  Every such point has an affinely independent positive support of
size at most five, and every smaller support extends to a five-column basis
of \(A\).
\end{lemma}

\begin{proof}
If \(p=\sum_t\alpha_t p^{(t)}\) is a convex combination of extreme laws and
\(q^{(t)}\) simultaneously dominates \(p^{(t)}\), then
\(q=\sum_t\alpha_tq^{(t)}\) has the same mass, first moments, and pair
moments as every \(q^{(t)}\).  Hence \(q\in\DPI(x)\), and linearity gives
\eqref{eq:simultaneous}.

If the columns \(a(S)\) in the positive support of an extreme \(p\) were
linearly dependent, a nonzero supported vector \(h\) would satisfy \(Ah=0\).
For sufficiently small \(\varepsilon>0\), both \(p+\varepsilon h\) and
\(p-\varepsilon h\) would be distinct feasible laws, contradicting
extremality.  Thus the support columns are independent and there are at most
five.  The basis-extension theorem extends any independent subset of columns
to a five-column basis.  Assigning zero weights to the added masks realizes
the original law on a face of that basis simplex.
\end{proof}

There are exactly \(3{,}008\) such five-column bases; this is independently
recomputed by exact rank enumeration.  All objects are equivariant under a
simultaneous permutation of the four ground elements.  We may therefore
relabel \(x,p\), apply the construction, and pull the resulting \(q\) back.
Hence it is enough to consider
\[
0\le x_1\le x_2\le x_3\le x_4\le1.
\tag{4.1}\label{eq:sorted}
\]
Ties cause no difficulty because every region below is closed.

\subsection{Two analytic regions}\label{subsec:analytic-regions}

\subsubsection{Low total marginal}

\begin{lemma}[Low-total region]\label{lem:low-total}
If \(s=\sum_i x_i\le1\), then \eqref{eq:CG4} holds.
\end{lemma}

\begin{proof}
Assume \eqref{eq:sorted}.  Prescribe candidate intersection moments by
\(m_i=x_i\), \(m_{ij}=x_ix_j\), and, for \(|T|\ge3\), let \(m_T\) be the
product of the two smallest marginals indexed by \(T\).  M\"obius inversion
gives the following potentially nonzero atom masses; all unlisted masses are
zero:
\[
\begin{array}{c@{\qquad}l}
\toprule
T&q_T\\
\midrule
\varnothing&1-s+x_4(s-x_4)\\
\{i\},\ i=1,2,3&x_i(1-x_4)\\
\{4\}&x_4(1-s+x_4)+x_3(x_1+x_2)\\
\{1,4\},\{2,4\}&x_1(x_4-x_3),\quad x_2(x_4-x_3)\\
\{3,4\}&x_3(x_4-x_1-x_2)+x_1x_2\\
\{1,3,4\},\{2,3,4\}&x_1(x_3-x_2),\quad x_2(x_3-x_1)\\
E&x_1x_2\\
\bottomrule
\end{array}
\tag{5.1}\label{eq:low-atoms}
\]
The empty mass is nonnegative because \(s\le1\).  All other masses are
manifestly nonnegative except possibly \(q_{\{3,4\}}\), and
\[
q_{\{3,4\}}
\ge x_3(x_3-x_1-x_2)+x_1x_2
=(x_3-x_1)(x_3-x_2)\ge0.
\]
The displayed nonnegative masses define a law \(Q\); equivalently, their
M\"obius transform is the prescribed moment vector and in particular has
total mass one.  The defining intersection moments show exactly that \(Q\) has marginals
\(x_i\) and pair intersections \(x_ix_j\), so it is pairwise independent.

For every \(A\subseteq E\),
\[
\Pr(Q\cap A\ne\varnothing)
\ge\frac34\sum_{i\in A}x_i.
\tag{5.2}\label{eq:hit-bound}
\]
For two coordinates with sum \(r\le1\), this follows from
\(4uv\le(u+v)^2\le u+v=r\).  For three coordinates with sorted marginals
\(a\le b\le c\), inclusion--exclusion and the chosen triple moment give
\[
\Pr(Q\cap A\ne\varnothing)=c+(a+b)-c(a+b),
\]
and the same inequality applies because \(c+a+b\le1\).  For \(A=E\), the
union probability is \(s-x_4(s-x_4)\), and the identical argument uses
\(c=x_4\), \(a+b=s-x_4\).  Singletons are immediate.

Normalize \(f(\varnothing)=0\), put \(a_i=f(\{i\})\), and define
\(h(\varnothing)=0\) and \(h(T)=\max_{i\in T}a_i\) for nonempty \(T\).
Monotonicity gives \(f\ge h\).  Layer cake and \eqref{eq:hit-bound} give
\[
\begin{aligned}
\E h(Q)
&=\int_0^\infty
\Pr\!\left(Q\cap\{i:a_i\ge u\}\ne\varnothing\right)\,du\\
&\ge\frac34\sum_i x_i a_i.
\end{aligned}
\tag{5.3}\label{eq:layer-cake}
\]
Normalized submodularity implies subadditivity and hence
\(f(T)\le\sum_{i\in T}a_i\).  Thus every law with marginals \(x\) has
expected value at most \(\sum_i x_i a_i\).  The law with mass \(x_i\) on
\(\{i\}\) and mass \(1-s\) on \(\varnothing\) attains the bound, so
\[
f^+(x)=\sum_i x_i a_i.
\]
Together with \eqref{eq:layer-cake}, this proves
\(f^{\mathrm{PI}}(x)\ge3f^+(x)/4\).  The argument uses no division and
includes all zero and deterministic boundary cases.
\end{proof}

\subsubsection{A large marginal}

We use the exact three-element theorem
\[
4u^{\mathrm{PI}}(y)\ge3u^+(y)
\tag{N3}\label{eq:n3}
\]
for every normalized nonnegative monotone submodular \(u\) and every real
\(y\in[0,1]^3\); see Ramachandra and Natarajan~\cite{RN2025} and the exact
campaign reconstruction.

\begin{lemma}[Large-coordinate region]\label{lem:large-coordinate}
If \(x_i=t\ge2/3\) for some coordinate, then \eqref{eq:CG4} holds.
\end{lemma}

\begin{proof}
Normalize \(f(\varnothing)=0\), put \(a=f(\{i\})\), and for
\(T\subseteq E\setminus\{i\}\) define
\[
G(T)=(1-t)f(T)+t f(T\cup\{i\}).
\tag{5.4}\label{eq:G}
\]
Here \(T\mapsto f(T)\) is deletion of coordinate \(i\), while
\(T\mapsto f(T\cup\{i\})\) is contraction.  Deletion and contraction preserve nonnegativity, monotonicity, and
submodularity; so does their convex combination.  Also
\(G(\varnothing)=ta\), and monotonicity gives
\(G-G(\varnothing)\ge0\).  Applying \eqref{eq:n3} to
\(G-G(\varnothing)\) yields the strengthened nonnormalized form
\[
4G^{\mathrm{PI}}(x_{-i})\ge3G^+(x_{-i})+ta.
\tag{5.5}\label{eq:G-n3}
\]
Adjoin to any three-coordinate PI law an independent Bernoulli-\(t\) bit at
coordinate \(i\).  This is a four-coordinate PI law and its expected
\(f\)-value is the expected \(G\)-value.  Hence
\[
f^{\mathrm{PI}}(x)\ge G^{\mathrm{PI}}(x_{-i}).
\tag{5.6}\label{eq:adjoin}
\]
Let \(P\) attain \(f^+(x)\), write \(X_i\) for its \(i\)-th bit and
\(Y=S\setminus\{i\}\), and set
\[
d(Y)=f(Y\cup\{i\})-f(Y).
\]
Diminishing returns and monotonicity give \(0\le d(Y)\le a\).  If only
\(X_i\) is independently resampled with mean \(t\), its expected value is
\[
R=\E[(1-t)f(Y)+t f(Y\cup\{i\})]
=\E G(Y)\le G^+(x_{-i}).
\tag{5.7}\label{eq:R}
\]
Moreover,
\[
\begin{aligned}
f^+(x)-R
&=\E[(X_i-t)d(Y)]\\
&\le(1-t)\E[X_i d(Y)]\\
&\le t(1-t)a.
\end{aligned}
\tag{5.8}\label{eq:cov-loss}
\]
Combining \eqref{eq:G-n3}--\eqref{eq:cov-loss},
\[
4f^{\mathrm{PI}}(x)
\ge3f^+(x)+ta(3t-2)
\ge3f^+(x).
\]
This includes \(t=2/3\), \(t=1\), and \(a=0\).
\end{proof}

By Lemmas~\ref{lem:low-total} and~\ref{lem:large-coordinate}, after sorting
it remains only to cover the closed central region
\[
0\le x_1\le x_2\le x_3\le x_4\le\frac23,
\qquad
\sum_i x_i\ge1.
\tag{C}\label{eq:central}
\]
Using a closed superset is deliberate: the two analytic regions already
cover the equality walls, and overlap creates no gap.

\subsection{The finite normalized-function test}
\label{subsec:cone}

For \(i\notin S\), write \(S+i=S\cup\{i\}\); for distinct
\(i,j\notin S\), write \(S+i+j=S\cup\{i,j\}\).  Also write
\(\Delta_i g(S)=g(S\cup\{i\})-g(S)\).  The normalized cone \(\Cfour\) is
defined exactly by the four terminal marginal inequalities
\[
g(E)-g(E\setminus\{i\})\ge0
\tag{M}\label{eq:terminal-monotonicity}
\]
and the \(24\) elementary diminishing-returns inequalities
\[
g(S+i)+g(S+j)-g(S)-g(S+i+j)\ge0
\tag{DR}\label{eq:DR}
\]
for \(i<j\) and \(S\subseteq E\setminus\{i,j\}\).  Indeed,
\eqref{eq:DR} telescopes to all diminishing-return inequalities and hence all
submodularity inequalities; together with terminal monotonicity it also
telescopes to every nonnegative marginal.

\begin{proposition}[Finite cone certificate]\label{prop:cone}
The pointed \(15\)-dimensional cone cut out by
\eqref{eq:terminal-monotonicity} and \eqref{eq:DR} has exactly \(41\)
primitive extreme rays in \(11\) \(S_4\)-orbits.
\end{proposition}

\noindent
\textit{Exact finite premise.}
The standard-library double-description transcript reconstructs the \(28\)
rows, starts from a certified simplicial cone, inserts the remaining \(13\)
halfspaces, and ends with \(41\) rays.  Rank-\(14\) integer minors certify
every final ray and every facet.  The cone-certificate SHA-256 is
\[
\texttt{4d99cba4266f6a9afbae04a9fff7a411cf450dfd78fd261d222a6d3a9d5ba13d},
\]
and the primitive-ray SHA-256 is
\[
\texttt{9f0f3d8e71de639468ea5915a4adfdff4056321040c7221def886e201624eff1}.
\]
The verifier reconstructs the defining rows and replays each
double-description transition; the certificate is therefore a completeness
certificate, not merely a supplied feasible ray list.

For fixed \(p,q\), conic generation makes \eqref{eq:simultaneous}
equivalent to its \(41\) ray inequalities, with the same \(q\) in every
inequality.  Ten rays are the coverage functions
\[
h_A(S)=\one[S\cap A\ne\varnothing],
\qquad 1\le|A|\le2.
\tag{6.1}\label{eq:or-rays}
\]
Their inequalities are automatic for every \(p\in\D(x)\) and every
\(q\in\DPI(x)\).  For \(A=\{i\}\), the slack is \(x_i\).  For
\(A=\{i,j\}\), write \(u=x_i\), \(v=x_j\), and \(s=u+v\).  Then
\[
\E_qh_A=s-uv,
\qquad
\E_ph_A\le\min(1,s).
\]
If \(s\le1\), then \(uv\le s^2/4\le s/4\).  If \(s\ge1\), put
\(a=1-u\), \(b=1-v\); then \(a+b\le1\), \(ab\le1/4\), and
\(s-uv=1-ab\ge3/4\).  Thus only \(31\) ray inequalities remain for the
central certificate.

\subsection{The exact central certificate}\label{subsec:central-certificate}

\begin{proposition}[Central certificate; computed exactly]
\label{prop:central-certificate}
For every five-column basis \(B\), every \(\lambda\) in its closed standard
simplex whose induced marginal vector lies in \eqref{eq:central}, and the
corresponding unrestricted law \(p\), at least one of the eight recipes below
is a law \(q\in\DPI(x)\) whose eleven atom masses and all \(31\)
nonautomatic simultaneous ray slacks are nonnegative.
\end{proposition}

The remainder of this subsection specifies the finite object and the exact
checks establishing Proposition~\ref{prop:central-certificate}.

Encode subsets by masks \(0,\ldots,15\), with bit \(i-1\) representing
element \(i\).  Fix a five-mask first-moment basis
\(B=(B_0,\ldots,B_4)\) and barycentric weights
\[
\lambda_k\ge0,
\qquad
\sum_{k=0}^4\lambda_k=1,
\qquad
p_{B_k}=\lambda_k.
\tag{7.1}\label{eq:lambda}
\]
Then
\[
x_i=\sum_k\lambda_k\one_{i\in B_k}.
\tag{7.2}\label{eq:x-lambda}
\]

\subsubsection{Eight explicit PI recipes}

Let
\[
m(x)=(1,x_1,x_2,x_3,x_4,x_1x_2,x_1x_3,x_1x_4,
x_2x_3,x_2x_4,x_3x_4)^T.
\tag{7.3}\label{eq:moment-vector}
\]
For an \(11\)-mask set \(R\), let \(M_R\) be the \(11\times11\) matrix whose
mask column records mass, four singleton indicators, and six pair
indicators.  The eight recipe supports are
\begin{verbatim}
R_0 = (0,3,4,6,8,9,10,12,13,14,15)
R_1 = (1,4,6,7,8,10,11,12,13,14,15)
R_2 = (1,2,4,7,8,9,10,12,13,14,15)
R_3 = (0,1,2,4,8,9,10,12,13,14,15)
R_4 = (0,1,2,3,4,5,6,9,10,12,15)
R_5 = (1,2,3,4,5,8,9,10,12,14,15)
R_6 = (1,3,5,6,7,8,10,11,12,13,15)
R_7 = (0,1,2,4,5,8,9,10,12,14,15)
\end{verbatim}
Every \(M_R\) is exactly invertible.  Recipe \(R\) is the law supported on
\(R\) with atom vector
\[
q_R(x)=M_R^{-1}m(x).
\tag{7.4}\label{eq:recipe}
\]
If its eleven entries are nonnegative, equations
\eqref{eq:moment-vector}--\eqref{eq:recipe} make it a probability law with
the required first and pair moments.  For each of the \(31\) nonautomatic
cone rays \(r\), impose
\[
4\sum_{S\in R}q_{R,S}(x)r(S)
-3\sum_{k=0}^4\lambda_kr(B_k)\ge0.
\tag{7.5}\label{eq:ray-slack}
\]
Thus a recipe has \(42\) constraints: eleven atom inequalities and \(31\)
dominance inequalities.  Since \(x\) is affine in \(\lambda\), every
constraint is a rational polynomial of degree at most two in \(\lambda\).

\subsubsection{Closed chamber cover}

Inside the standard \(\lambda\)-simplex impose the homogeneous linear forms
equivalent to \eqref{eq:central}.  Split first on
\[
T=x_1+x_2+x_3-1.
\]
The cell \(T\le0\) is one chamber.  On \(T\ge0\), split independently on
\[
P=x_1+x_2-1,
\qquad
H=1+x_1-x_2-x_3.
\]
This gives four more chambers.  Because both signs are non-strict, these five
closed chambers cover the whole central cell, including all walls.

Two basis-specific refinements are used, always with both closed signs:
\begin{verbatim}
B=(1,6,8,10,12):       2 lambda_0-lambda_1-lambda_3 = 0;

B=(2,5,8,9,12), in the T>=0, P<=0, H>=0 chamber:
    lambda_0+lambda_1-lambda_2+lambda_3-lambda_4 = 0,
   -lambda_0+lambda_1-lambda_2-lambda_3+lambda_4 = 0.
\end{verbatim}
The latter two forms equal, respectively,
\(2(x_1+x_2-1/2)\) and \(2(x_3-1/2)\).  Both-sign splitting preserves
coverage even on their common zero sheets.

For each basis and sign chamber, exact active-set enumeration lists every
vertex of the bounded rational polytope.  A point is a vertex exactly when
four active inequality normals, together with \(\sum\lambda=1\), have full
rank; hence enumerating all four-row choices is complete.  Every nonempty
bounded cell has a vertex.  Exact pulling triangulation then covers the cell:
recursively choose the least vertex and cone it over pulling triangulations
of all facets not containing it.  Facets are recovered from the original
inequalities by exact affine-rank tests.

\subsubsection{Exact quadratic checks on a simplex}

Let \(v_0,\ldots,v_d\) be a simplex and write
\(\lambda=\sum_i y_iv_i\), where \(y_i\ge0\) and \(\sum_i y_i=1\).  A
quadratic constraint \(h\) has the exact representation
\[
h(\lambda)=y^THy.
\tag{7.6}\label{eq:quadratic}
\]
The diagonal controls are \(h(v_i)\), and the off-diagonal controls are
\[
H_{ij}=2h\!\left(\frac{v_i+v_j}{2}\right)
-\frac{h(v_i)+h(v_j)}2.
\tag{7.7}\label{eq:bernstein}
\]
These are the degree-two Bernstein controls;
\eqref{eq:bernstein} is an identity, not a positivity heuristic.

\begin{lemma}[Quadratic simplex lemma]\label{lem:quadratic-simplex}
Exact enumeration of every face-stationary affine set decides whether
\eqref{eq:quadratic} is nonnegative on the whole simplex.
\end{lemma}

\begin{proof}
If \(h\) is negative, its compact-simplex minimum is negative.  On the
relative interior of its support face \(I\), Lagrange stationarity gives
\[
(Hy)_i=(Hy)_j\qquad(i,j\in I),
\qquad
\sum_{i\in I}y_i=1.
\tag{7.8}\label{eq:stationarity}
\]
The exact checker intersects \eqref{eq:stationarity} with \(y\ge0\) by
enumerating basic feasible solutions.  If \(y,z\) satisfy the same stationary
system, then \(Hy=\alpha\mathbf 1\) and \(Hz=\beta\mathbf 1\) on \(I\), so
symmetry gives
\[
\alpha=z^THy=y^THz=\beta.
\]
Consequently \(y^THy=\alpha=\beta=z^THz\): the quadratic value is constant
on the stationary affine set.  Testing one exact feasible basic point per
nonempty stationary set therefore detects a negative minimum if and only if
one exists.  This also handles singular stationary systems and boundary
minima.
\end{proof}

\subsubsection{Accepted simplex labels}

For every simplex the certificate records one of the following exact labels.
\begin{enumerate}
\item \textbf{Ordinary recipe.}  All \(42\) constraints of one recipe are
nonnegative on the simplex by Lemma~\ref{lem:quadratic-simplex}.

\item \textbf{Complementary switch.}  Recipe \(R\) has at most one
uncertified constraint \(u\), recipe \(R'\) has at most one uncertified
constraint \(v\), and \(v=-cu\) identically for an exact rational \(c>0\).
If \(u\ge0\), the first recipe works; if \(u\le0\), the second works.  At
\(u=0\), both blockers vanish.  The recurring global identity is
\[
q_{R_7,\varnothing}=-2q_{R_5,\{1,2\}},
\]
where mask \(3\) represents \(\{1,2\}\).

\item \textbf{Farkas/disjunctive label.}  For each selected recipe, all
constraints except a listed blocker set are nonnegative.  For every selection
of one blocker \(h_j\) from each selected recipe, the certificate gives exact
weights \(c_j\ge0\), not all zero, such that
\[
\sum_jc_jh_j\ge0
\tag{7.9}\label{eq:farkas}
\]
throughout the simplex, again by
Lemma~\ref{lem:quadratic-simplex}.  If all selected recipes failed at one
point, choose a strictly negative blocker from each.  The left side of
\eqref{eq:farkas} would be strictly negative, a contradiction.  The
one-blocker-per-recipe special case is called a Farkas label.
\end{enumerate}
The replay verifier recomputes all non-blocker constraints; it does not trust
exploratory safety flags.

\subsubsection{Retained exact certificate}

The retained certificate covers all \(3{,}008\) bases and has the following
exact totals:
\begin{center}
\begin{tabular}{lr}
\toprule
closed chamber/sign records&15,048\\
nonempty records&9,288\\
empty records&5,760\\
initial simplices/leaves&17,933\\
ordinary-recipe leaves&17,906\\
switch leaves&21\\
Farkas/disjunctive leaves&6\\
maximum subdivision depth&0\\
\bottomrule
\end{tabular}
\end{center}
The canonical payload is
\path{verification/central/central-cell-cover-v1.json.gz}.  Its sizes and
hashes are
\begin{center}
\scriptsize
\begin{tabular}{@{}ll@{}}
\toprule
raw JSON bytes&1,352,318\\
gzip bytes&115,720\\
raw JSON SHA-256&\texttt{c18f2096019a80ec495209863bdd8e46faa8e9e812c11e014cfc334b2a02f7ec}\\
gzip SHA-256&\texttt{dbad1fcaee0e0285b448e1179c73ea747b5225c52dbd60da70064a1cc6727958}\\
\bottomrule
\end{tabular}
\end{center}
The exporting scan used exact \texttt{Fraction} arithmetic.  A standalone
standard-library Python bundle reconstructed and replayed all four ranges,
including \(286{,}253\) label-local exact quadratic checks.  A separate
Node.js implementation using only built-ins and custom \texttt{BigInt}
rationals independently reconstructed the \(3{,}008\) bases, every chamber
H-polytope, all vertices and pulling simplices, the eight recipe inverses and
\(42\) constraints, and all label semantics.  It performed \(754{,}395\)
exact face-stationary quadratic tests and found no discrepancy.  Thus the
simplex argument covers all real barycentric points, not merely rational or
sampled points.

\subsection{Assembly of the proof}\label{subsec:assembly}

\begin{proof}[Proof of Theorem~\ref{thm:main}]
First normalize \(f\) using Lemma~\ref{lem:normalization}.  Fix a real
\(x\in[0,1]^4\).  If \(\sum_i x_i\le1\),
Lemma~\ref{lem:low-total} proves \eqref{eq:CG4}.  If
\(\max_i x_i\ge2/3\), Lemma~\ref{lem:large-coordinate} proves it.

Otherwise fix an extreme \(p\in\D(x)\), sort the coordinates, and relabel
\(p\) accordingly.  Then \(x\) belongs to the closed central region
\eqref{eq:central}.  By Lemma~\ref{lem:extreme-laws}, the support of \(p\)
extends with zero weights to one of the \(3{,}008\) five-column bases.  Its
barycentric vector lies in at least one of the closed chamber/sign cells and
then in at least one simplex of that cell's pulling triangulation.  The
simplex label supplies at least one recipe whose eleven atom probabilities
and \(31\) nonautomatic dominance slacks are all nonnegative.  The ten
remaining ray slacks are automatic.  Thus one
\(q\in\DPI(x)\) satisfies all \(41\) ray inequalities and hence
simultaneously dominates \(p\) on all of \(\Cfour\).

Convexly combine these \(q\)'s to handle an arbitrary \(p\in\D(x)\).
Lemma~\ref{lem:minimax} then gives \eqref{eq:CG4} for every normalized
function.  Lemma~\ref{lem:normalization} restores arbitrary
\(f(\varnothing)\ge0\), and undoing the coordinate permutation restores the
original labeling.  This covers all real marginals, zero barycentric weights,
chamber walls, the equality cases \(\sum x_i=1\) and
\(\max x_i=2/3\), and all ties among sorted coordinates.

For sharpness, take the two-element OR function embedded on four elements and
\[
x=(1/2,1/2,0,0).
\]
An unrestricted law may choose exactly one of the first
two elements, so \(f^+(x)=1\).  Pairwise independence fixes their intersection
at \(1/4\), so their union has probability \(3/4\) and
\(f^{\mathrm{PI}}(x)=3/4\).  Hence the ratio is \(4/3\).
\end{proof}

\subsection{Scope and residual trust}\label{subsec:trust}

The argument is computer-assisted.  Its finite premises use exact rational
arithmetic rather than floating-point evidence, and the decisive central
certificate has a second-language replay.  The residual trust base consists
of the elementary polyhedral and quadratic lemmas stated above, the retained
certificate bytes, JSON/gzip parsing, and the correctness of the
standard-language runtimes and exact-arithmetic implementations.  This is not
a Lean, Coq, or Isabelle formalization.  The appropriate status is
\emph{candidate resolution pending expert review}.

\ifdefined\CGFullPaper
\else
\begin{thebibliography}{9}
\bibitem{RN2025}
A.~Ramachandra and K.~Natarajan,
\emph{Pairwise independent correlation gap},
Operations Research Letters \textbf{60} (2025), Article 107255.
\href{https://doi.org/10.1016/j.orl.2025.107255}{doi:10.1016/j.orl.2025.107255}.
\end{thebibliography}
\fi
\ProofOnlyEndDocument
